% !TEX root = ../supplement_0421052099.tex
% ---------------------------------------------------------------------
% APPENDIX C. Created 2026-09-23 with the other two, under a hard
% budget of 20-22 pages for abstract-to-conclusion. RELOCATED from
% Section 4: the two baseline implementations at the width a
% reimplementation needs, and the definitions behind the groundedness
% tiers. Nothing was dropped.
%
% WHAT STAYED IN sec:baselines: the claim each baseline licenses, which
% is what Sections 5-8 reason from. What moved here is how each was
% built. Same rule as Appendix B.
%
% THE KAPPA SENTENCES DID NOT MOVE, and must not: check_paper_numbers
% takes section_body("sec:protocol") and requires \kappa = 0.6995, the
% word "misses" against $0.70$, and the shortfall 0.0005 to be inside
% sec:protocol itself. Four ways to lose that were recorded when the
% block was written; moving the sentence to an appendix would be a
% fifth.
%
% IT LIVES IN sections/ ON PURPOSE -- see appendix-instrument.tex.
% ---------------------------------------------------------------------

\section{Protocol detail}\label{app:protocol}

\Cref{sec:setup} states what each baseline licenses and how the
comparison is controlled. This appendix records how the two non-trivial
baselines were built and how groundedness was measured.

\subsection{The static graph baseline}\label{app:graphrag}

Static GraphRAG~\cite{edge2024local} fetches a fixed-radius
neighbourhood around the linked topic entities in one shot, reranks it
by embedding similarity to the question, and answers from the top $30$
facts. Two properties of that implementation bound what it licenses, and
we state them rather than leave a reader to infer them from its scores.
The radius is \emph{one logical hop}: mediators are traversed
transparently, so a mediator path costs two graph edges while remaining
one hop, and a genuine two-hop chain is outside its reach by
construction. And expansion retrieves at most $100$ edges per topic
entity with no ordering criterion, so on a high-degree entity the
retained edges are an arbitrary sample rather than the best or the
first; at least one topic entity is truncated on $72.5\%$ of the $800$
test questions. The baseline therefore licenses one narrow claim, the
one \Cref{sec:baselines} draws, and not a claim about what a two-hop
static retriever would score. That experiment was not run.

\subsection{The agentic baseline}\label{app:tog}

Think-on-Graph~\cite{sun2024think} performs model-guided beam search
over relations and entities, reimplemented on this environment's tools
rather than quoted. Beam width and depth follow the original paper, $3$
and $3$. At each depth it makes one relation-pruning call per frontier
entity, one entity-pruning call per selected relation, and one
sufficiency call, so a full depth at width $3$ costs up to thirteen
calls, and the $25$-call cap of \ref{sec:budgets} admits one full depth
and part of a second. Two widths beneath the beam are this
reimplementation's own rather than the original's: it offers the pruning
prompt the first $40$ relations of an entity and the first $20$
neighbours of a relation, against the $300$ and $200$ AGR sees
(\ref{sec:budgets}). Reimplementing is a stronger comparison than citing
the original figures, because graph, model, budget, and scoring rule are
held constant and only the algorithm varies. It also means its numbers
are not comparable to the published ones, and we do not present them as
such. All baselines were debugged on development data and certified
before any test question ran, against a pre-specified diagnostic
targeting the one failure that would be an implementation defect rather
than a paradigm limit: retrieval found the gold answer and the system
hedged anyway.

\subsection{How groundedness is measured}\label{app:groundedness}

\textbf{Groundedness} is measured at two levels. The structural tier is
deterministic and applies to every asserted entity in all $4{,}000$
records: an entity is grounded if and only if it exists as a node in the
environment and connects to a topic entity within four hops. The
semantic tier asks the harder question of whether the entity is
supported \emph{as the answer}. For a stratified sample of $60$ answers
per system per dataset, it retrieves the graph's own shortest connecting
paths and has a language model judge whether any path expresses the
relationship the question asks about, under the governing instruction
that mere connectedness is not support. Because path retrieval takes
only the three shortest paths, semantic-tier rates are conservative
lower bounds. They are depressed uniformly across systems and read as a
between-system comparison rather than as absolute support rates.
\Cref{sec:protocol} reports the judge's validation against a human
annotator, including the threshold it missed.
